\section{Différents algorithmes de détection}

\subsection{Principe général du CFAR (Constant False Alarm Rate)}

\subsubsection{Estimation du bruit}

Le niveau du clutter (fouillis radar : échoss parasites) est estimé par la moyenne arithmétique des échantillons dans la fenêtre de référence donnée par : 
\begin{equation}
{Q}_{\text{clutter}} = \frac{1}{M} \sum_{i=1}^{M} x_i
\end{equation}

\begin{description}
  \item[${Q}_{\text{clutter}}$ :] estimation du niveau moyen du clutter (ou du bruit de fond)
  \item[$N$ :] nombre total de cellules de référence
  \item[$x_i$ :] puissance mesurée dans la $i^{\text{ème}}$ cellule de référence
\end{description}

Si on considère la détection quadratique, la FDP (fonction de densité de probabilité) de la variable aléatoire à la sortie du détecteur, X, suit
la loi Exponentielle donnée par :
\begin{equation}
f_X(x) = \frac{1}{2\sigma^2} e^{-x / 2\sigma^2}, \quad x \ge 0
\end{equation}




\subsubsection{Calcul du seuil de détection}
\subsubsection{Probabilité de fausse alarme}
